\section{Seleção cronológica de obras, periódicos e páginas inéditos}

A quarta rodada amplia o levantamento a partir de três endereços indicados pelo autor: a coleção geral da Brasiliana/USP, uma captura da página de entrada do arquivo de comunidades do Orkut e o portal da Wayback Machine. Adotou-se uma seleção cronológica de títulos ainda não examinados, incluindo sites, fóruns e blogs, conforme delimitação expressa do autor. A unidade editorial é uma referência por obra, periódico ou site. Essa regra difere da quota por acervo e década da rodada anterior, cujo estado de conclusão permanece publicado. Não se reapresentam títulos anteriores como novidades por terem sido reencontrados em outro repositório.

\subsection{Seleção, exclusão de repetições e cronologia}

O inventário anterior de obras e jornais, incluindo as indicações auditadas sem ocorrência confirmada, foi conservado como lista de exclusão. Os títulos selecionados foram confrontados com esse inventário; volumes de uma mesma obra e números de um periódico não foram multiplicados para aumentar o total. A busca dirigida por \forma{riso} na Brasiliana retornou 1.078 itens na consulta de descoberta registrada. Resultados ordenados por data foram usados para descoberta, com inspeção de exemplares escolhidos; não se leu integralmente esse conjunto. O mesmo identificador de coleção fornecido pelo autor foi acessado no endereço atual \url{https://digital.bbm.usp.br/handle/bbm/1}. A falha de certificado no endereço antigo não foi resolvida mediante desativação da verificação de segurança.

A seleção executada reúne dezessete títulos: quatorze da Brasiliana, um blog, um fórum e uma comunidade pública do Orkut. O conjunto não constitui um contínuo anual: o intervalo dos testemunhos selecionados vai de 1645 a 2017, e a data da consulta é 2026. O exemplar de 1645 é o mais antigo confirmado nesta seleção, não a primeira ocorrência portuguesa de \forma{riso} nem o livro mais antigo da biblioteca. A existência de materiais quinhentistas no catálogo não permite atribuir-lhes a ocorrência sem recuperar o trecho.

As referências complementares distribuem-se em treze menções lexicais, dois registros lexicográficos e duas grafias miméticas. As duas últimas integram o núcleo de representações; sua inclusão antecedeu a ampliação anual. No balanço consolidado, o núcleo contém 121 registros: 23 E1, 95 E2 e três E3. As demais referências documentam usos, sentidos ou tematizações do riso; não são convertidas em exemplos de \forma{kkkk}.

\subsection{Séculos XVII e XVIII: leitura das fontes e história da escrita}

Na \textit{Chronica da Companhia de Iesv na provincia de Portvgal}, volume I, de Baltasar Teles, edição de 1645, a página 266 do PDF contém uma referência a risos e zombarias. A leitura foi conferida na imagem. O OCR também oferece um falso candidato em outra página, lendo como \forma{risos} uma sequência que não autentica o fenômeno. No trecho selecionado, a tipografia emprega o chamado \emph{s longo}, reconhecido mecanicamente como \forma{f}. A normalização para \forma{risos} é uma decisão de transcrição, não uma alegação de mudança fonética de \forma{rifos} para \forma{risos} \cite{r4chronicacompanhiaiesvportugal}.

O volume VII do \textit{Vocabulario Portuguez e latino}, de Rafael Bluteau, de 1720, apresenta o verbete \forma{Riso} na página impressa 345, correspondente à página 353 do PDF. O verbete aborda sinais corporais e manifestações da alegria, mas suas explicações devem ser lidas como história da lexicografia e das ideias de sua época, sem reproduzir sua fisiologia como ciência atual. No \textit{Diccionario da lingua portugueza} reformado por Antônio de Moraes Silva, de 1789, o volume II contém o verbete \forma{RISO}, na página 354 do PDF. São obras lexicográficas relacionadas, não duas provas independentes de uma mesma genealogia digital \cite{r4vocabularioportuguezlatino,r4diccionariolinguaportuguezamorais}.

Em \textit{Viola de Lereno}, volume I, de 1798, a expressão \forma{Mostra em riso moderado} ocorre na página 259 do PDF. A passagem descreve o riso e a exposição dos dentes; não grafa uma risada por repetição de sílabas. A seleção mantém apenas uma referência para a obra, em vez de somar seus volumes e reedições como títulos novos \cite{r4violadelereno}.

\subsection{Periódicos, teatro, poesia e prosa dos séculos XIX e XX}

O recorte inclui três periódicos inéditos no inventário anterior: \textit{O Patriota}, no conjunto de números de 1813; \textit{Gazeta Pernambucana}, de 1822; e \textit{Nitheroy}, de 1836. As expressões selecionadas, respectivamente, remetem a ser digno de riso, a merecer risada e a um riso de alegria. São testemunhos lexicais localizados nas páginas 25, 26 e 140 dos PDFs correspondentes. A \textit{Gazeta Pernambucana} não foi confundida com o \textit{Diario de Pernambuco} nem com a \textit{Gazeta do Rio de Janeiro} \cite{r4opatriota,r4gazetapernambucana,r4nitheroy}.

Nas \textit{Reflexões dramaticas para uso dos canditados que se dedicam a scena}, de João Caetano, edição de 1837, a recomendação de evitar rir às gargalhadas em determinado papel aparece na página 22 do PDF, página impressa 16. Trata-se de reflexão sobre representação cênica, não de proibição gramatical geral. \textit{Primeiros cantos}, de Gonçalves Dias, edição de 1846, oferece uma referência lexical na página impressa 47, PDF 51. \textit{Gennesco}, de 1866, \textit{A fallencia}, de 1902, e \textit{Negrinha}, de 1920, acrescentam usos em prosa com referências a risos, gargalhadas e reações de personagens. As imagens dos trechos foram inspecionadas, sem considerar o texto de OCR uma transcrição automaticamente fidedigna \cite{r4reflexoesdramaticas,r4primeiroscantos,r4gennesco,r4afallencia,r4negrinha}.

Em \textit{Poesia, 1958/1978}, de Ana Hatherly, edição de 1980, a página 35 do PDF contém \forma{O riso tem lágrimas}. O intervalo do título não permite atribuir esse poema a um ano específico entre 1958 e 1978 sem outro cotejo. O registro foi datado pelo exemplar consultado. Em \textit{Rubens Borba e Moraes: Anotações de um Bibliófilo}, de Cristina Antunes, publicação de 2017, a página impressa 144, PDF 146, reproduz uma referência a poesias que serviam de riso. A redação remete a material antigo, mas o testemunho consultado é a publicação de 2017; não se retrodata automaticamente a frase a um manuscrito anterior \cite{r4poesia19581978,r4rubensborbamoraesanotacoesbibliofilo}.

\subsection{Wayback, blog e fórum: três datas que não devem ser confundidas}

Uma página preservada pode oferecer três datas diferentes: a data que o conteúdo exibe, a data de criação do espaço que o hospeda e a data da captura. Além disso, solicitar um instante à Wayback pode levar a outro registro disponível. O manifesto conserva tanto a URL solicitada quanto a URL efetivamente servida. Essa diferença foi decisiva na triagem de uma página incorporada do \textit{Humor Tadela}: uma consulta iniciada numa captura de dezembro de 1998 conduziu a um recurso de novembro de 2002, sem ocorrência pertinente autenticada. O recurso não foi contado como risada documentada em 1998.

No \textit{Kibe Loco}, a postagem \textit{Uma chance (parte 2)} exibe a data de 27 de junho de 2007 e a menção \forma{não contendo o riso}. O conteúdo foi recuperado numa captura efetiva de 1º de julho de 2007, embora o pedido inicial ao endereço de entrada estivesse associado a 10 de junho. A referência comprova a presença lexical naquela postagem e não a publicação em junho de conteúdo que a própria página data posteriormente \cite{r4kibeloco}.

No fórum Cifra Club, a captura de 15 de setembro de 2010 do tópico \textit{como afinar meio tom abaixo} preserva uma resposta com \forma{kkkkkkkkkkkkk}, identificada na interface como enviada em 12 de julho de 2010, às 15h21. O tópico foi iniciado em novembro de 2003. A data de abertura não pode ser transferida à resposta acrescentada em 2010. O registro fornece texto primário eletrônico preservado de escrita com \forma{k} repetido; não antecipa para 2003 a ocorrência mais antiga deste levantamento \cite{r4forumcifraclub}.

\subsection{Orkut: página de entrada, comunidade e descrição preservada}

A captura indicada pelo autor, de 1º de outubro de 2014, apresenta a entrada do arquivo de comunidades do Orkut. A página descreve o período de funcionamento do serviço, de janeiro de 2004 a setembro de 2014. Esse texto institucional não é uma conversa de 2004 e não comprova sozinho uma forma específica de riso \cite{r4orkutindex}.

A navegação pelo índice levou à comunidade pública \textit{Kiko Dança Psy-TRANCE}. Sua descrição contém \forma{hehehe} na captura efetiva de 4 de outubro de 2014. Embora a página indique criação da comunidade em 24 de fevereiro de 2006, a descrição pode ter sido editada posteriormente. Na ausência de uma versão anterior do trecho, o testemunho foi datado pela captura de 2014, estabelecendo apenas que a sequência estava preservada até esse momento. Não se atribui a redação a 2006, nem se apresenta o exemplo como prova de uma primeira utilização no Orkut \cite{r4orkutcomunidades}.

Os três espaços digitais cumprem papéis diferentes nesta seleção: referência lexical em blog, resposta de fórum com marcador de riso e descrição de comunidade com grafia mimética. Sua comparação esclarece funções e condições de preservação, mas não reconstrói uma cadeia causal entre plataformas. Uma captura arquivada tampouco representa todas as interações que o serviço abrigou.

\subsection{Reprodução e limites da ampliação}

O repositório \repo\ publica a lista de exclusão anterior, a tabela de títulos inéditos e o resumo da rodada em \nolinkurl{revision4_exclusions.csv}, \nolinkurl{revision4_new_titles.csv} e \nolinkurl{revision4_summary.json}. As aquisições têm localizadores e SHA-256 no manifesto geral. A lista cronológica aparece no apêndice. A seleção não é um censo da Brasiliana ou da web preservada, não é equivalente à seleção por quota de acervo e década e não permite estimar prevalências. Não se aplicou uma fórmula amostral para justificar dezessete títulos: o número é o resultado desta seleção dirigida, condicionada pela novidade em relação ao inventário, pela diversidade temporal e pela autenticação dos trechos.
